\subsection{Foundational Definitions of a System}

The mathematical formalization of systems theory emerges from diverse epistemological traditions, each seeking to articulate the fundamental nature of systemic relationships through rigorous set-theoretic frameworks. We present here the principal definitions that constitute the conceptual architecture of General Systems Theory.

\subsubsection{The Verbal-Conceptual Definition}

Mesarovic \citep{Mesarovic1972} articulates the primordial intuition underlying all systems thinking:

\begin{quote}
\emph{A system is (or reflects the existence of) a relationship between the objects (of study).}
\end{quote}

This formulation reveals the essential phenomenological stance: systems are not \emph{things-in-themselves}, but rather patterns of interrelation discerned through observational praxis. The definition achieves generality precisely through its abstraction from particular physical instantiations, concerning itself instead with the structural configurations of relationships.

Mesarovic further refines this characterization to emphasize its teleological dimension:

\begin{quote}
\emph{A system is a relationship among objects described (or, specified, defined) in terms of information-processing and decision-making concepts.}
\end{quote}

\subsubsection{The Set-Theoretic Formalization}

The transition from verbal description to mathematical precision constitutes the foundational gesture of modern systems theory. Mesarovic \citep{Mesarovic1972} proposes the canonical set-theoretic definition:

\begin{align*}
S \subseteq V_1 \times \cdots \times V_n
\end{align*}

where $\times$ denotes the Cartesian product. The coordinate sets $V_1, \ldots, V_n$ are designated as \emph{system objects}. This formulation encapsulates several profound insights:

\begin{itemize}
\item An object $V_i$ represents the totality of all possible appearances of an attribute of the phenomenon under consideration.
\item The system $S$ specifies which combinations of object appearances can occur simultaneously when objects are interrelated to form the given system.
\item The relation $S$ is not merely an arbitrary subset, but rather embodies the structural constraints that govern the co-variation of system attributes.
\end{itemize}

Mesarovic and Takahara further develop this framework, emphasizing that systems are defined on the set-theoretic level with minimal mathematical structure \citep{Mesarovic1975}. They write:

\begin{quote}
A family of sets is given, $\mathcal{V} = \{V_i : i \in I\}$ where $I$ is the index set, and a system, defined on $\mathcal{V}$, is a proper subset of $\times \mathcal{V}$:
\begin{align*}
S \subseteq \times \{V_i : i \in I\}
\end{align*}
\end{quote}

\subsubsection{The Input-Output System}

For systems exhibiting causal structure, the formalism specializes to an input-output configuration. The most widely studied case considers a system consisting of two objects \citep{Mesarovic1975}:

\begin{align*}
S \subseteq X \times Y
\end{align*}

where $X$ represents the input object (stimuli) and $Y$ represents the output object (responses). This partitioning introduces the fundamental distinction between cause and effect, transforming the abstract relational structure into a framework amenable to dynamical analysis.

When $S$ is a function, i.e., with every input $x \in X$ there is associated a unique output, the system is designated a \emph{functional system}:

\begin{align*}
S: X \to Y
\end{align*}

This functional characterization permits the introduction of \emph{internal mechanisms}---constructive specifications that determine output generation for any given input through differential equations, transformation tables, or algorithmic procedures.

\subsubsection{The Multirelation Approach}

Yi Lin \citep{Lin1999} introduces an alternative formalization emphasizing the multiplicity of relations within a system. A system is defined as an ordered pair:

\begin{align*}
S = (M, R)
\end{align*}

where $M$ is a set of \emph{objects} and $R$ is a set of \emph{relations} on $M$. For each relation $r \in R$, there exists an ordinal number $n = n(r)$ such that $r \subseteq M^n$, where $M^n$ denotes the $n$-fold Cartesian product of $M$.

This formulation explicitly acknowledges that systems typically exhibit multiple simultaneous relationships of varying arities, capturing the richness of interconnection patterns that characterize complex phenomena.

\subsubsection{Bunge's Compositional Definition}

Bunge introduces environmental considerations into the system definition \citep[as cited in][]{Lin1999,Bunge1979}. Given a nonempty set $T$, an ordered triple $W = (C, E, S)$ constitutes a system over $T$ if and only if:

\begin{itemize}
\item $C$ and $E$ are mutually disjoint subsets of $T$ (composition and environment, respectively)
\item $S$ is a nonempty set of relations on $C \cup E$
\end{itemize}

This tripartite structure distinguishes internal composition from external environment, while acknowledging that systemic relations may span both domains.